% GENERATED FILE. Edit canonical lessons, then run npm run generate:books.
% canonical-source-hash: fnv1a64:c21436c36c570474
% canonical-lessons: ML-C344-ticket-venam, ML-C344-saukaryam, ML-W344-au-sign, ML-C344-bus-kittiyilla, ML-R344-at-the-station

\chapter{Tickets, Seats and Trouble}
\label{ch:ml-tickets-seats-and-trouble}
\hlchaptermodality{344}

\begin{chapteropening}
\textbf{By the end of this chapter:} \emph{I can buy a ticket and ask about times and seats, and report a missed bus, a late train or a lost bag and ask for help.}

Complete ML-R344-at-the-station: buy a ticket and ask when the train leaves, report a missed bus and a lost bag, ask for help with -āmō, and read the new sign \ml{ൗ}.
\end{chapteropening}

\section[\texorpdfstring{\ml{രണ്ട്} \ml{ടിക്കറ്റ്} \ml{വേണം} (raṇṭŭ ṭikkaṟṟŭ vēṇaṁ)}{രണ്ട് ടിക്കറ്റ് വേണം (raṇṭŭ ṭikkaṟṟŭ vēṇaṁ)}]{\ml{രണ്ട്} \ml{ടിക്കറ്റ്} \ml{വേണം} (raṇṭŭ ṭikkaṟṟŭ vēṇaṁ) — two tickets, please}

\label{lesson:ML-C344-ticket-venam}



\begin{quote}
Say \textquotedblleft{}a ticket\textquotedblright{}, \textquotedblleft{}a seat\textquotedblright{}, and \textquotedblleft{}at what time?\textquotedblright{} (\emph{ṭikkaṟṟŭ}, \emph{sīṟṟŭ}, \emph{etra maṇikkŭ}.)
\end{quote}

\begin{grammarlens}[title={Grammar Lens: at the counter}]
Name the place with \textbf{-\ml{ലേക്ക്}}, \textquotedblleft{}to\textquotedblright{}, then how many, then \textbf{\ml{വേണം}}:

\begin{quote}
\textbf{\ml{കൊച്ചിയിലേക്ക്} \ml{രണ്ട്} \ml{ടിക്കറ്റ്} \ml{വേണം}.}
\end{quote}

\begin{quote}
\emph{kocciyilēkkŭ raṇṭŭ ṭikkaṟṟŭ vēṇaṁ.} — \textquotedblleft{}Two tickets to Kochi, please.\textquotedblright{}
\end{quote}

For the way back, use \textbf{\ml{മടങ്ങുക}}, \textquotedblleft{}to return\textquotedblright{}, in its \textbf{-\ml{ആൻ}} form, with \textbf{-\ml{ഉം}}, \textquotedblleft{}too\textquotedblright{}. At a counter you will also hear the English \emph{return ticket}.

\begin{quote}
\textbf{\ml{മടങ്ങാനും} \ml{ടിക്കറ്റ്} \ml{വേണം}.}
\end{quote}

\begin{quote}
\emph{maṭaṅṅānuṁ ṭikkaṟṟŭ vēṇaṁ.} — \textquotedblleft{}I need a ticket to come back too.\textquotedblright{}
\end{quote}

Then the price: \textbf{\ml{എത്രയാണ്}?} (\emph{etrayāṇŭ?}), \textquotedblleft{}how much is it?\textquotedblright{}
\end{grammarlens}

\begin{grammarlens}[title={Grammar Lens: times and seats}]
\begin{quote}
\textbf{\ml{ട്രെയിൻ} \ml{എത്ര} \ml{മണിക്ക്} \ml{പുറപ്പെടും}?}
\end{quote}

\begin{quote}
\emph{ṭreyin etra maṇikkŭ puṟappeṭuṁ?} — \textquotedblleft{}What time does the train leave?\textquotedblright{}
\end{quote}

\textbf{\ml{പുറപ്പെടുക}} is \textquotedblleft{}to set out\textquotedblright{}; the future \textbf{\ml{പുറപ്പെടും}} asks about the timetable. So \textbf{\ml{എപ്പോൾ} \ml{എത്തും}?} (\emph{eppōḷ ettuṁ?}), \textquotedblleft{}when will it arrive?\textquotedblright{}, and \textbf{\ml{സീറ്റ്} \ml{ഉണ്ടോ}?} (\emph{sīṟṟŭ uṇṭō?}), \textquotedblleft{}is there a seat?\textquotedblright{}
\end{grammarlens}

\subsection*{Your turn}
Say these aloud:

\begin{itemize}
\raggedright
  \item \textquotedblleft{}Two tickets to Kochi, please.\textquotedblright{}
  \item \textquotedblleft{}I need a ticket to come back too. How much is it?\textquotedblright{}
  \item \textquotedblleft{}What time does the train leave? When will it arrive?\textquotedblright{}
  \item \textquotedblleft{}Is there a seat?\textquotedblright{}
\end{itemize}

\begin{tcolorbox}[breakable,colback=teal!4,colframe=teal!35!black,title={Before you move on}]
\textquotedblleft{}To Kochi\textquotedblright{}? (\textbf{\ml{കൊച്ചിയിലേക്ക്}}, \emph{kocciyilēkkŭ}, with \textbf{-\ml{ലേക്ക്}}.) \textquotedblleft{}A ticket to come back too\textquotedblright{}? (\emph{maṭaṅṅānuṁ ṭikkaṟṟŭ vēṇaṁ.}) \textquotedblleft{}What time does the train leave?\textquotedblright{} (\emph{\textbf{ṭreyin etra maṇikkŭ puṟappeṭuṁ?}})
\end{tcolorbox}
\section[\texorpdfstring{\ml{സൗകര്യം} (saukaryaṁ)}{സൗകര്യം (saukaryaṁ)}]{\ml{സൗകര്യം} (saukaryaṁ) — convenience, facilities}

\label{lesson:ML-C344-saukaryam}



\begin{quote}
Before the new one: say \textquotedblleft{}two tickets to Kochi, please\textquotedblright{}, then \textquotedblleft{}what time does the train leave?\textquotedblright{}
\end{quote}

\subsection*{You'll want to know: \ml{സൗകര്യം}}
\textbf{\ml{സൗകര്യം}} — \emph{saukaryaṁ} — \textquotedblleft{}convenience\textquotedblright{}, and in the plural \textquotedblleft{}facilities\textquotedblright{}.

\begin{quote}
\textbf{\ml{സ്റ്റേഷനിൽ} \ml{നല്ല} \ml{സൗകര്യങ്ങൾ} \ml{ഉണ്ട്}.}
\end{quote}

\begin{quote}
\emph{sṟṟēṣanil nalla saukaryaṅṅaḷ uṇṭŭ.} — \textquotedblleft{}The station has good facilities.\textquotedblright{}
\end{quote}

With the dative it says what suits you:

\begin{quote}
\textbf{\ml{നാളെ} \ml{എനിക്ക്} \ml{സൗകര്യമാണ്}.}
\end{quote}

\begin{quote}
\emph{nāḷe enikkŭ saukaryamāṇŭ.} — \textquotedblleft{}Tomorrow suits me.\textquotedblright{}
\end{quote}

\textbf{\ml{നിങ്ങളുടെ} \ml{സൗകര്യം} \ml{പോലെ}}, \textquotedblleft{}as suits you\textquotedblright{}, is how a polite Malayali leaves the choice to you.

Its first syllable carries a vowel sign you have not written alone. The next lesson takes it out.

\begin{cousinweb}
\textbf{\ml{സൗകര്യം}} (\emph{saukaryaṁ}) $\leftarrow$ Sanskrit \emph{saukarya}, \textquotedblleft{}ease, facility\textquotedblright{}, a tatsama with Malayalam's noun ending. Monier-Williams derives \emph{saukarya} from \emph{su-kara}, \textquotedblleft{}easy to do\textquotedblright{}: \emph{su-}, \textquotedblleft{}well, easily\textquotedblright{}, and \emph{kara}, \textquotedblleft{}doing\textquotedblright{}.
\end{cousinweb}

\subsection*{Your turn}
Say these aloud:

\begin{itemize}
\raggedright
  \item \emph{saukaryaṁ}, with \emph{au} as in \textquotedblleft{}cow\textquotedblright{}
  \item \textquotedblleft{}The station has good facilities.\textquotedblright{}
  \item \textquotedblleft{}Tomorrow suits me\textquotedblright{}, then \textquotedblleft{}as suits you\textquotedblright{}
\end{itemize}

\begin{tcolorbox}[breakable,colback=teal!4,colframe=teal!35!black,title={Before you move on}]
What does \textbf{\ml{സൗകര്യം}} mean? (\textbf{Convenience; facilities.}) \textquotedblleft{}Tomorrow suits me\textquotedblright{}? (\emph{\textbf{nāḷe enikkŭ saukaryamāṇŭ.}})
\end{tcolorbox}
\section[\texorpdfstring{\ml{ൗ} (au)}{ൗ (au)}]{\ml{ൗ} (au) — the sign inside \ml{സൗകര്യം}}

\label{lesson:ML-W344-au-sign}



\begin{quote}
Before the new one: write \textbf{\ml{െ}} and \textbf{\ml{ൈ}} once each, and name them. (The \emph{e} sign and the \emph{ai} sign, both printed before their consonant.)

Now say \textbf{\ml{സൗകര്യം}} (\emph{saukaryaṁ}), \textquotedblleft{}convenience\textquotedblright{}. This lesson takes the sign in its first syllable and puts it in your hand.
\end{quote}

\subsection*{Script you'll notice: \ml{ൗ}}
\textbf{\ml{ൗ}} — the vowel sign \emph{au}. It is printed after its consonant: \textbf{\ml{സ}} + \textbf{\ml{ൗ}} $\to$ \textbf{\ml{സൗ}}, \emph{sau}.

Its formal name is the \emph{au} length mark. Older books also put \textbf{\ml{െ}} in front of the consonant for the same \emph{au}, so you may meet both in print; they are read alike. The mark alone is the modern spelling.

You have written this shape already, as the tail of \textbf{\ml{ഈ}} and \textbf{\ml{ഔ}}: a small lobe, a notch, and a broader lobe sweeping back along its foot. Here it stands on its own.

\subsection*{Writing: \ml{ൗ}}
One recorded hand writes \textbf{\ml{ൗ}} in one stroke, with no lift: start low on the left and climb over the small lobe, dip into the notch, climb over the broader lobe and down its right side, and sweep back left along the foot.

Write \textbf{\ml{ൗ}} yourself, slowly, and larger than it is printed. Copy the printed shape and keep the order above.

\subsection*{Your turn}
\begin{itemize}
\raggedright
  \item \emph{Find:} \ml{ൗ} inside \textbf{\ml{സൗകര്യം}}
  \item \emph{Trace it:} \ml{ൗ} three times, saying \emph{au} as you finish each one
  \item \emph{Write it:} \textbf{\ml{സൗ}}, then \textbf{\ml{കൗ}}, then \textbf{\ml{സൗകര്യം}}
\end{itemize}

\begin{tcolorbox}[breakable,colback=teal!4,colframe=teal!35!black,title={Before you move on}]
Which sign is this — \ml{ൗ}? (\textbf{au}, printed after its consonant.) Which word did it come from? (\textbf{\ml{സൗകര്യം}}, \emph{saukaryaṁ}.) Which letters end in the same tail? (\textbf{\ml{ഈ}}, \textbf{\ml{ഊ}} and \textbf{\ml{ഔ}}.)
\end{tcolorbox}
\section[\texorpdfstring{\ml{എനിക്ക്} \ml{ബസ്} \ml{കിട്ടിയില്ല} (enikkŭ basŭ …)}{എനിക്ക് ബസ് കിട്ടിയില്ല (enikkŭ basŭ …)}]{\ml{എനിക്ക്} \ml{ബസ്} \ml{കിട്ടിയില്ല} (enikkŭ basŭ kiṭṭiyilla) — I missed the bus}

\label{lesson:ML-C344-bus-kittiyilla}



\begin{quote}
Say \textquotedblleft{}to get\textquotedblright{}, \textquotedblleft{}to lose\textquotedblright{}, and \textquotedblleft{}it's late\textquotedblright{}. (\emph{kiṭṭuka}, \emph{naṣṭappeṭuka}, \emph{vaiki}.) Then the cry when something goes wrong. (\emph{ayyō!})
\end{quote}

\begin{grammarlens}[title={Grammar Lens: it did not come to me}]
Malayalam usually turns \textquotedblleft{}I missed the bus\textquotedblright{} round: the bus did not come to you, and you stand in the dative:

\begin{quote}
\textbf{\ml{എനിക്ക്} \ml{ബസ്} \ml{കിട്ടിയില്ല}.}
\end{quote}

\begin{quote}
\emph{enikkŭ basŭ kiṭṭiyilla.} — \textquotedblleft{}I missed the bus.\textquotedblright{}
\end{quote}

Or the bus went: \textbf{\ml{ബസ്} \ml{പോയി}} (\emph{basŭ pōyi}), \textquotedblleft{}the bus has gone\textquotedblright{}. Delays and cancellations use verbs you know:

\begin{quote}
\textbf{\ml{ട്രെയിൻ} \ml{ഒരു} \ml{മണിക്കൂർ} \ml{വൈകി}.}
\end{quote}

\begin{quote}
\emph{ṭreyin oru maṇikkūr vaiki.} — \textquotedblleft{}The train was an hour late.\textquotedblright{}
\end{quote}

\textbf{\ml{എന്റെ} \ml{ടിക്കറ്റ്} \ml{റദ്ദാക്കി}} (\emph{enṟe ṭikkaṟṟŭ ṟaddākki}), \textquotedblleft{}my ticket was cancelled\textquotedblright{}: no one is named, as in English.
\end{grammarlens}

\begin{grammarlens}[title={Grammar Lens: lost and broken, and help}]
\begin{quote}
\textbf{\ml{അയ്യോ}, \ml{എന്റെ} \ml{ബാഗ്} \ml{നഷ്ടപ്പെട്ടുപോയി}.}
\end{quote}

\begin{quote}
\emph{ayyō, enṟe bāgŭ naṣṭappeṭṭupōyi.} — \textquotedblleft{}Oh no, my bag has got lost.\textquotedblright{}
\end{quote}

\textbf{\ml{പോയി}} after a past form says it went and is gone, and you are sorry about it: \textbf{\ml{ഫോൺ} \ml{കേടായിപ്പോയി}}, \textquotedblleft{}the phone has broken down\textquotedblright{}.

To ask for help, take the \textbf{-\ml{ആം}} of \textbf{\ml{വിളിക്കാം}}, \textquotedblleft{}I'll call\textquotedblright{}, and add \textbf{-\ml{ഓ}}: \textquotedblleft{}could you …?\textquotedblright{}

\begin{quote}
\textbf{\ml{ക്ഷമിക്കണം}, \ml{എന്നെ} \ml{സഹായിക്കാമോ}?}
\end{quote}

\begin{quote}
\emph{kṣamikkaṇaṁ, enne sahāyikkāmō?} — \textquotedblleft{}Excuse me, could you help me?\textquotedblright{}
\end{quote}
\end{grammarlens}

\subsection*{Your turn}
Say these aloud:

\begin{itemize}
\raggedright
  \item \textquotedblleft{}I missed the bus\textquotedblright{}, then \textquotedblleft{}the bus has gone\textquotedblright{}
  \item \textquotedblleft{}The train was an hour late. My ticket was cancelled.\textquotedblright{}
  \item \textquotedblleft{}Oh no, my bag has got lost.\textquotedblright{}
  \item \textquotedblleft{}Excuse me, could you help me?\textquotedblright{}, then \textquotedblleft{}could you call?\textquotedblright{} — \emph{viḷikkāmō?}
\end{itemize}

\begin{tcolorbox}[breakable,colback=teal!4,colframe=teal!35!black,title={Before you move on}]
\textquotedblleft{}I missed the bus\textquotedblright{}? (\emph{\textbf{enikkŭ basŭ kiṭṭiyilla.}}) What does \textbf{\ml{പോയി}} add after \emph{naṣṭappeṭṭu}? (\textbf{it is gone, and you are sorry.}) \textquotedblleft{}Could you help me?\textquotedblright{} (\emph{\textbf{enne sahāyikkāmō?}})
\end{tcolorbox}
\section[(review)]{Review — at the station}

\label{lesson:ML-R344-at-the-station}



\begin{quote}
Say \textquotedblleft{}convenience; facilities\textquotedblright{}. (\emph{saukaryaṁ}.) Write \textbf{\ml{ൗ}} once from memory, then \textbf{\ml{സൗ}}.
\end{quote}

\subsection*{Your turn}
\begin{itemize}
\raggedright
  \item \emph{Say it:} \textquotedblleft{}Two tickets to Kochi, please. How much is it?\textquotedblright{}
  \item \emph{Say it:} \textquotedblleft{}What time does the train leave? Is there a seat?\textquotedblright{}
  \item \emph{Say it:} \textquotedblleft{}I missed the bus. The train was an hour late.\textquotedblright{}
  \item \emph{Say it:} \textquotedblleft{}Oh no, my bag has got lost. Excuse me, could you help me?\textquotedblright{}
  \item \emph{Say it:} \textquotedblleft{}Tomorrow suits me.\textquotedblright{}
  \item \emph{Say it:} \textquotedblleft{}Where is the station? How do I get there?\textquotedblright{}
  \item \emph{Write it:} \textbf{\ml{സൗകര്യം}}, then \textbf{\ml{റോഡ്}}
\end{itemize}

\begin{tcolorbox}[breakable,colback=teal!4,colframe=teal!35!black,title={Before you move on}]
\textquotedblleft{}Two tickets to Kochi\textquotedblright{}? (\emph{kocciyilēkkŭ raṇṭŭ ṭikkaṟṟŭ vēṇaṁ.}) \textquotedblleft{}I missed the bus\textquotedblright{}? (\emph{\textbf{enikkŭ basŭ kiṭṭiyilla.}}) \textquotedblleft{}Could you help me?\textquotedblright{} (\emph{\textbf{enne sahāyikkāmō?}}) \textquotedblleft{}Facilities\textquotedblright{}? (\textbf{\ml{സൗകര്യങ്ങൾ}}, with the sign \textbf{\ml{ൗ}}.)
\end{tcolorbox}
